\section{Weekly Summary}

\begin{tcolorbox}[colback=yellow!10,colframe=black!50,title=AI-Generated Content]
\noindent The summary below was written by Claude (Anthropic) from that week's regression outputs and series descriptions. It has not been reviewed or fact-checked by a human, and should be read as an automated, informal recap -- not analysis.
\end{tcolorbox}

Another glorious week of shoving two complete strangers into a regression and asking them to explain each other's life choices. As always, the official disclaimer up front: these pairings are chosen by the cold, uncaring hand of a random number generator, nobody here is claiming anything causes anything, and if you walk away believing Danish solar stocks control Polish recessions, that is a you problem, not a statistics problem. Spurious is the house style. Let us proceed.

Monday kicked things off with the OMX Copenhagen Alternative Energy Price Index versus a discontinued Polish recession indicator, which is already the kind of blind date that ends with both parties checking their phones. The raw regression coughed up a technically-significant p-value attached to an R-squared of roughly 0.06, which is the quantitative equivalent of "we have literally nothing in common but sure, we kissed." Amazingly, the detrended version actually held up and even nudged the R-squared up to 0.07, so this one survived the trend-stripping gauntlet. I would love to tell you this means something. It means two series explain seven percent of each other. Hold your applause. Tuesday was mercifully honest: car-part export prices versus Chilean purchases of U.S. agency bonds produced a p-value of 0.43 raw and a majestic 0.90 detrended, which is the number you get when two series refuse to acknowledge each other's existence. Beautiful. No notes.

Then came the midweek festival of trend mirages, our bread and butter. Wednesday paired consumer motor-vehicle leases with foreign-held mortgage-backed securities and posted a show-stopping raw R-squared of 0.81 -- gasp! -- only for the detrended version to deflate to 0.44. Still significant, which is suspicious enough to be mildly interesting, but that thirty-seven-point faceplant is pure "they both went up over time" energy. Thursday did the same trick harder: Covered Employment versus a Midwest house price index strutted in at 0.79 raw, then collapsed to 0.17 once we took away its shared-drift crutch, which is the statistical arc of a soufflé. Friday was the purest mirage of the week -- child poverty in Leon County, Florida versus high-tier San Francisco home prices, a gaudy 0.42 raw that detrended all the way down to 0.006 with a p-value of 0.71. That is not a relationship collapsing; that is a relationship that was never there filing for divorce.

Saturday gave us the week's most entertaining betrayal: Mid-Atlantic private employment versus a management-compensation cost index, a towering 0.77 raw R-squared that evaporated to 0.009 detrended, p-value 0.44. Two series that both trend upward because the economy tends to do that, caught red-handed. And Sunday closed us out with retail sales percent-change versus San Angelo, Texas home prices, which never even bothered to fake it -- 0.005 raw, 0.004 detrended, both profoundly insignificant. A fittingly humble finish. So the week's scoreboard: a parade of raw "hits" that crumbled the instant we asked them to prove themselves, and Monday's weird little copper-colored survivor that technically held up while explaining almost nothing. Treat all of it as the random noise it proudly is, and I will see you next week for more of the same.
